\chapter{메뉴: \texorpdfstring{\code{menu}}{menu}}\label{ch:rust-menu}

\code{menu} 모듈은 봄베를 돌리기 전에 해독자가 종이 위에서 하던 일을 한다. 크립을 놓을 자리를 추리고,
크립과 암호문의 글자 관계를 그래프로 만들고, 기계에 걸 수 있는 크기로 다듬는다.

\section{크립을 놓을 자리: \texorpdfstring{\code{possible\_offsets}}{possible\_offsets}}

\begin{lstlisting}[language=Rust,caption={\code{menu.rs}: 크립이 놓일 수 있는 위치.}]
pub fn possible_offsets(ciphertext: &[u8], crib: &[u8]) -> Vec<usize> {
    if crib.is_empty() || crib.len() > ciphertext.len() {
        return Vec::new();
    }
    (0..=ciphertext.len() - crib.len())
        .filter(|&o| crib.iter().zip(&ciphertext[o..]).all(|(p, c)| p != c))
        .collect()
}
\end{lstlisting}

읽는 법: 가능한 모든 시작 위치 \code{o}에 대해, 크립의 글자 \code{p}와 그 아래 놓이는 암호문 글자 \code{c}를
짝지어(\code{zip}) \emph{모든} 짝이 서로 다르면(\code{all(p != c)}) 그 위치를 남긴다.

\Why 에니그마는 어떤 글자도 자기 자신으로 암호화하지 않는다(앞 장). 그러므로 크립의 글자와 암호문의 글자가
같은 자리가 하나라도 있으면 그 위치는 있을 수 없다.

\General 낱말 맞추기 놀이에서 ``이 자리에는 이 글자가 올 수 없다''는 단서로 후보를 지워 나가는 것과 같다.
단서 하나는 약하지만 크립이 길수록 단서가 많아진다. 틀린 위치가 우연히 살아남을 확률은 글자 하나마다
$25/26$씩 곱해진다.

\Ours 예제 전문(121글자)에서 16글자 크립 \texttt{WETTERVORHERSAGE}를 놓을 수 있는 위치는 106곳이고,
이 검사로 46곳이 지워져 60곳이 남는다.

\section{메뉴의 자료 구조}

\begin{lstlisting}[language=Rust,caption={\code{menu.rs}: 선과 메뉴(주석 일부 생략).}]
pub struct Edge {
    pub plain: u8,
    pub cipher: u8,
    pub pos: u16,
}

pub struct Menu {
    pub offset: usize,
    pub crib: Vec<u8>,
    pub edges: Vec<Edge>,
    pub test_letter: u8,
    pub input_wire: u8,
}
\end{lstlisting}

\code{Edge}는 메뉴의 선 하나다. 크립의 \code{pos}번째 자리(0부터)에서 평문 글자 \code{plain}이 암호문 글자
\code{cipher}가 되었다는 사실이다. 봄베에서는 선 하나에 스크램블러 하나가 걸린다. \code{Menu}는 선들의
목록에 두 가지를 더한 것이다. 전류를 넣고 관찰할 글자 \code{test\_letter}와, 그 글자의 케이블 26가닥
가운데 전압을 걸 선 \code{input\_wire}다.

\code{Menu::from\_crib}은 크립의 글자마다 \code{Edge}를 하나씩 만든다. 그러고 나서 테스트 글자를 고른다.
가장 큰 덩어리 안에서 선이 가장 많이 모이는 글자다. 실제 운용에서도 선이 가장 많이 모이는 글자를
골랐다~\cite[7]{us6812report1944}. 전류가 여러 방향으로 한꺼번에 퍼져 나가 모순이 빨리 드러나기 때문이다.

\section{덩어리와 루프}

메뉴의 품질을 재는 수가 둘 있다. 서로 이어진 글자들의 \term{덩어리}{connected component} 수와 루프 수다.

\begin{lstlisting}[language=Rust,caption={\code{menu.rs}: 루프의 수.}]
pub fn loops(&self) -> usize {
    let letters = self.letters().count_ones() as usize;
    (self.edges.len() + self.components().len()).saturating_sub(letters)
}
\end{lstlisting}

\Why 루프의 수가 $(\text{선 수}) - (\text{글자 수}) + (\text{덩어리 수})$인 이유는 이렇다. 글자 $n$개를 한
덩어리로 잇는 데 필요한 선은 최소 $n-1$개이고, 그때 루프는 없다(나무 모양). 거기에 선을 하나 더하면 이미
이어져 있던 두 글자 사이에 두 번째 길이 생기므로 루프가 정확히 하나 생긴다. 그러므로 루프의 수는
``나무에 필요한 것보다 남는 선의 수''다(그림~\ref{fig:g-loops}).

\begin{figure}[ht]
\centering
\begin{tikzpicture}[x=1cm,y=1cm]
  \foreach \dx/\ttl/\extra in {0/{글자 4, 선 3: 루프 0}/0, 5.3/{글자 4, 선 4: 루프 1}/1, 10.6/{글자 4, 선 5: 루프 2}/2} {
    \begin{scope}[xshift=\dx cm]
      \node[ltr] (a) at (0,0) {A}; \node[ltr] (b) at (1.6,1.0) {B};
      \node[ltr] (c) at (3.2,0) {C}; \node[ltr] (d) at (1.6,-1.0) {D};
      \draw (a) -- (b) -- (c); \draw (a) -- (d);
      \ifnum\extra>0 \draw[cblue,line width=1.3pt] (d) -- (c); \fi
      \ifnum\extra>1 \draw[cblue,line width=1.3pt] (b) -- (d); \fi
      \node[note] at (1.6,-1.9) {\ttl};
    \end{scope}
  }
\end{tikzpicture}
\caption{선을 하나 더할 때마다 루프가 하나 는다. 파란 선이 나무에 더해진 선이다.}
\label{fig:g-loops}
\end{figure}

덩어리를 세는 \code{components}는 \term{서로소 집합}{union-find}이라는 표준적인 방법을 쓴다. 처음에는 글자마다
자기만의 묶음에 있고, 선 하나를 볼 때마다 그 양 끝 글자의 묶음을 합친다. 다 보고 나면 남은 묶음이 덩어리다.

\General 서로소 집합은 ``누가 누구와 같은 무리인가''를 묻는 문제에 두루 쓰인다. 친구 관계에서 무리를
나누거나, 통신망에서 어느 장비끼리 이어져 있는지를 찾는 일이 같은 꼴이다.

\Ours 루프가 중요한 이유는 봄베의 원리에 있다. 테스트 글자의 스테커 짝을 하나 가정하고 루프를 한 바퀴 돌면
다시 테스트 글자의 스테커 짝이 계산되어 나오는데, 이것이 처음 가정과 같아야 한다. 틀린 드럼 위치에서 우연히
같아질 확률은 약 $1/26$이다. 그래서 루프 하나가 거짓 정지를 약 $1/26$로 줄인다. 덩어리가 중요한 이유는
전류가 떨어진 덩어리로는 건너가지 못하기 때문이다. 테스트 글자가 속한 덩어리의 선만 쓸모가 있다.

\section{기계에 맞게 다듬기: \texorpdfstring{\code{select}}{select}}

봄베의 스크램블러 수는 한정되어 있다(영국 봄베는 뱅크당 12개). \code{select(max\_edges)}는 선의 수를 제한한
부분 메뉴를 고른다. 방법은 \term{탐욕법}{greedy}이다. 그 순간 가장 좋아 보이는 선을 하나씩 고른다.

\begin{lstlisting}[language=Rust,caption={\code{menu.rs}: \code{select}의 핵심(일부).}]
while chosen.len() < max_edges && !remaining.is_empty() {
    let inside = |e: &Edge| have >> e.plain & 1 == 1 && have >> e.cipher & 1 == 1;
    let touches = |e: &Edge| have >> e.plain & 1 == 1 || have >> e.cipher & 1 == 1;
    let pick = if let Some(i) = remaining.iter().position(inside) {
        Some(i)
    } else {
        remaining.iter().enumerate()
            .filter(|(_, e)| touches(e))
            .max_by_key(|(_, e)| {
                let new = if have >> e.plain & 1 == 1 { e.cipher } else { e.plain };
                (deg[new as usize], std::cmp::Reverse(e.pos))
            })
            .map(|(i, _)| i)
    };
    ...
}
\end{lstlisting}

\code{have}는 지금까지 고른 글자의 집합을 26비트 정수 하나로 나타낸 것이다. \code{have >> e.plain \& 1 == 1}은
``글자 \code{e.plain}이 집합에 들어 있는가''다. 글자의 집합을 비트로 다루는 이 표현은 다음 장에서 케이블의
통전 상태를 나타낼 때에도 쓰인다.

규칙은 둘이다. 양 끝이 \emph{둘 다} 이미 고른 글자인 선(\code{inside})이 있으면 먼저 고른다. 그런 선은
새 글자를 더하지 않고 루프만 하나 닫는다. 없으면 한쪽 끝만 닿는 선(\code{touches}) 가운데 새로 들어올 글자의
연결이 가장 많은 것을 고른다. 나중에 루프를 닫을 가능성이 크기 때문이다.

\Ours 예제에서 크립을 23글자(\texttt{WETTERVORHERSAGEFUERDIE})로 늘리고 \code{select(14)}를 하면
선 14개, 글자 11개, 루프 4개의 메뉴가 나온다. 16글자 크립에서 만든 메뉴(선 14개, 글자 14개, 루프 1개)와
스크램블러 수는 같지만, Enigma~I의 모든 로터 순서를 돌렸을 때 정지가 157회에서 2회로 줄어든다.
